\section{Equal-DOF baseline comparisons}
\label{sec:results}

\subsection{Comparison setup and scope}

Five cases use a common evaluator: the one-dimensional linear problem, the one-dimensional cubic
problem, the two-dimensional sinh problem, the one-dimensional pure \(p=4\) problem, and the
two-dimensional pure \(p=4\) problem. The comparison variable is the number of independent global
coefficients: retained atoms for CGA, ridge features held fixed for RFM, and independent
finite-element coefficients after the constant mode is removed where required. Equal degrees of
freedom (DOF) do not imply equal runtime, memory, score evaluations, or nonlinear solver work, and
no cross-method timing conclusion is drawn from runs performed in different environments.
The supplementary numerical data at
\url{https://github.com/RonzolYu/cga4PDE} give the prescribed realization count, the
available converged realizations, and the solver status for each width.

FEM uses continuous Lagrange elements on uniform one-dimensional meshes or structured triangular
two-dimensional meshes.  P1 and P3 provide low- and high-order mesh-based references in the main
figures. P2 results are retained in the supplementary material. P3 is the most accurate among the tested uniform-mesh FEM variants
at each largest common coefficient count reported below. The linear and cubic FEM curves are close
at high resolution, which is consistent with the common manufactured profile and an \(H^1\) error
dominated by the finite-element approximation.

For each case, the RFM uses a nested family of normalized \(\relu^3\) ridges and fits the outer
coefficients. Ten random-feature realizations are prescribed for each problem. Curves show the median and
interquartile range over available converged realizations at each width, and the supplementary numerical data give
the prescribed count together with solver status. All ten
realizations reach the displayed ranges for the linear, cubic, sinh, and one-dimensional pure
\(p=4\) problems. In the two-dimensional pure \(p=4\) problem, seven realizations are available at
\(N=128\) and four are available at \(N=256\); consequently its statistical curve ends at 256 and its terminal band is based
on four available realizations.  No distributional statement is made at \(N=512\).

The CGA and RFM comparison curves join the computed powers of two listed in the numerical setup and the
terminal retained state needed to define the largest common comparison width, while FEM uses its own mesh degrees of
freedom. The figures contain no common-grid interpolants, theoretical reference lines, or
convergence-order-versus-DOF curves. The reported panel range ends at the terminal CGA width fixed
by the stopping rule, so that endpoint is determined by the stopping rule and the available method
states rather than by the plotted error values. Each method is shown only over the states it
actually reached. The equal-DOF tables use linear interpolation of log(error) against log(DOF)
between adjacent computed values when an exact coefficient count is unavailable; no extrapolation
is used.

\begin{table}[!htbp]
\centering
\small
\caption{Training and evaluation parameters for the equal-DOF comparison.}
\label{tab:protocol}
\begin{tabularx}{\textwidth}{@{}p{0.26\textwidth}X@{}}
\toprule
Component & Setting \\
\midrule
RFM & Nested \(\relu^3\) prefixes; the discrete homogeneous normalization defined in \Cref{sec:numerics}; centered for pure \(p\); ten prespecified random seeds, with prescribed and available-realization counts reported. \\
FEM & Continuous P1/P2/P3; structured meshes; assembly order 10, increased to 80 for the one-dimensional multifrequency cases; relative/absolute Newton tolerances \(10^{-11}\)/\(10^{-12}\). \\
Training quadrature & One dimension: 256 cells, order 6; two dimensions: \(63\times63\) cells, tensor-product order 3. \\
Evaluation quadrature & One dimension: 2048 cells, order 8; two dimensions: \(64\times64\) cell partition, order 3. \\
Independent refinement & One dimension: 3072 cells, order 8; two dimensions: \(80\times80\) cell partition, order 3. \\
Common DOFs & \(8,16,32,64,128,256\), with 512 used only when every method required by the comparison is bracketed by computed states. \\
\bottomrule
\end{tabularx}
\end{table}
\FloatBarrier

\subsection{Unified evaluator, retained states, and fixed denominators}

The comparison states are re-evaluated after training by a frozen evaluator that is separate from
the quadrature used to construct the features or fit their coefficients.  The main evaluation rule
uses the quadrature listed in \Cref{tab:protocol}; an independent refinement is used for the independently evaluated quantities.  Thus, a solver-specific status is not interpreted as an independent accuracy bound, and a state is not removed merely because an independently evaluated quantity is unavailable.  The
set of retained states and the denominator conventions are summarized in
\Cref{tab:unified-audit-boundary,tab:rfm-fixed-denominators}.

The common evaluation retains 617 stored states, comprising 612 newly generated states and five archived
terminal states.  All 617 satisfy their solver-specific acceptance criterion and the independent coefficient-count
check.  The stricter signed energy/metric evaluation admits 611 states, while the common \(H^1\)-dual
residual threshold admits 394; these are different filters and are reported with their own
denominators.  The independent two-dimensional Bregman calculation is an energy-integration evaluation
on 138 states, not a replacement for the residual or optimizer checks.

For random-feature comparisons, the prescribed denominator is ten seeds at every configured width.
The displayed median and interquartile range are computed only from evaluated states that pass the
metric-specific checks, whereas missing rank-budget states remain failures in the \(10\)-seed
denominator.  In particular, the three C3 states missing at \(N=256\), the one C5 state missing at
\(N=128\), and the two C5 states missing at \(N=256\) are not dropped before forming success counts.
All other displayed RFM widths have \(10/10\) evaluated states; their strict residual counts can
still be smaller because residual qualification is a separate filter.

\begin{table}[!htbp]
\centering
\small
\caption{Unified re-evaluation boundaries and fixed denominators for the archived comparison states.}
\label{tab:unified-audit-boundary}
\begin{tabularx}{\textwidth}{@{}p{0.42\textwidth}X@{}}
\toprule
Audit item & Count or convention \\
\midrule
Retained states in the common audit & 617 total: 612 newly generated; 5 archived terminal \\
Native solver acceptance & (617/617); this records the native acceptance rule only \\
Independent coefficient-count check & (617/617) \\
Signed energy/metric audit & (611/617); six states remain excluded from the corresponding strict energy summary \\
Common H1-dual residual, threshold 2e-6 & (394/617); this is reported separately from native acceptance \\
Independent two-dimensional Bregman-energy audit & (137/138); this tests energy-integration stability and is not an optimization certificate \\
\bottomrule
\end{tabularx}
\end{table}

\begin{table}[!htbp]
\centering
\scriptsize
\caption{Prescribed random-feature denominators at widths where states are missing or further audit filtering is active.  A missing state remains in the denominator.}
\label{tab:rfm-fixed-denominators}
\begin{tabularx}{0.98\textwidth}{@{}lrrrr>{\raggedright\arraybackslash}X@{}}
\toprule
Case/width & Plan & Eval. & Missing & Energy & Strict H1 residual \\
\midrule
C3, (N=256) & 10 & 7 & 3 & (7/10) & (0/10) \\
C5, (N=128) & 10 & 9 & 1 & (9/10) & (0/10) \\
C5, (N=256) & 10 & 8 & 2 & (7/10) & (0/10) \\
\bottomrule
\end{tabularx}
\end{table}

\FloatBarrier

\subsection{Relative Sobolev errors}

For the linear, cubic, and sinh problems,
\Cref{fig:baseline-sobolev-1d,fig:baseline-sobolev-rest} report the relative \(H^1\) error. For the
two pure \(p=4\) problems, the corresponding panels report the relative
\(W^{1,4}\) error. These are Sobolev errors; the natural \(V\)-distance is reported separately below.

\begin{figure}[!htbp]
\centering
\begin{subfigure}[t]{0.49\textwidth}
\includegraphics[width=\linewidth]{C1_relative_h1_error.pdf}
\caption{Linear, \(d=1\).}
\end{subfigure}\hfill
\begin{subfigure}[t]{0.49\textwidth}
\includegraphics[width=\linewidth]{C2_relative_h1_error.pdf}
\caption{Cubic, \(d=1\).}
\end{subfigure}
\caption{Relative \(H^1\) error versus effective degrees of freedom for the one-dimensional
multifrequency problems. RFM curves show medians and interquartile bands over successful realizations
from ten prespecified runs. All markers are computed states.}
\label{fig:baseline-sobolev-1d}
\end{figure}

\begin{figure}[!htbp]
\centering
\begin{subfigure}[t]{0.32\textwidth}
\includegraphics[width=\linewidth]{C3_relative_h1_error.pdf}
\caption{Sinh, \(d=2\): relative \(H^1\) error.}
\end{subfigure}\hfill
\begin{subfigure}[t]{0.32\textwidth}
\includegraphics[width=\linewidth]{C4_relative_w1p_error.pdf}
\caption{Pure \(p=4\), \(d=1\): relative \(W^{1,4}\) error.}
\end{subfigure}\hfill
\begin{subfigure}[t]{0.32\textwidth}
\includegraphics[width=\linewidth]{C5_relative_w1p_error.pdf}
\caption{Pure \(p=4\), \(d=2\): relative \(W^{1,4}\) error.}
\end{subfigure}
\caption{Relative Sobolev errors for the two-dimensional sinh problem and the pure \(p=4\) problems.
The two-dimensional pure \(p=4\) RFM band is formed from the available realizations; its statistics at \(N=128\) and 256 use
seven and four available realizations, respectively.}
\label{fig:baseline-sobolev-rest}
\end{figure}
\FloatBarrier

At the largest common endpoint supported by the configured RFM widths and bracketing FEM states,
CGA has the smallest relative Sobolev error in the linear, cubic, sinh, and one-dimensional pure
\(p=4\) problems, for which all ten prescribed RFM solves succeeded.
The RFM-median/CGA ratios are 1.49, 1.52, 1.77, and 5.76, and the P3/CGA ratios are 4.24, 4.35, 5.64,
and 1.88. In the two-dimensional pure \(p=4\) problem the successful-run RFM median and the P3 value
are also above the CGA value, by factors
5.01 and 16.12, respectively, but the RFM statistic is based on only four of ten prescribed solves.
These are endpoint comparisons at equal coefficient count. The curves cross at smaller widths in
some cases, so the data do not support uniform dominance at every DOF.

\begin{table}[!htbp]
\centering\scriptsize
\caption{Relative Sobolev error at the largest common coefficient count reported for each problem. RFM statistics use successful realizations from ten prescribed runs and are reported as median [Q1,Q3]. The first four problems have ten successful runs; the two-dimensional pure \(p=4\) row is conditional on four successful runs and is not included in an aggregate ranking. Ratios larger than one favor CGA.}
\label{tab:baseline-terminal}
\begin{tabular}{@{}crrrrrr@{}}
\toprule
Problem & DOF & CGA & RFM median [Q1,Q3] & Success & RFM/CGA & FEM P3/CGA \\ \midrule
Linear, d=1 & 256 & $1.247\!\times\!10^{-4}$ & $1.853\!\times\!10^{-4}$ [$1.779\!\times\!10^{-4}$, $1.913\!\times\!10^{-4}$] & 10/10 & 1.49 & 4.24 \\
Cubic, d=1 & 256 & $1.215\!\times\!10^{-4}$ & $1.853\!\times\!10^{-4}$ [$1.779\!\times\!10^{-4}$, $1.913\!\times\!10^{-4}$] & 10/10 & 1.52 & 4.35 \\
Sinh, d=2 & 256 & $4.074\!\times\!10^{-3}$ & $7.195\!\times\!10^{-3}$ [$6.472\!\times\!10^{-3}$, $8.663\!\times\!10^{-3}$] & 10/10 & 1.77 & 5.64 \\
Pure p=4, d=1 & 128 & $1.579\!\times\!10^{-5}$ & $9.103\!\times\!10^{-5}$ [$6.108\!\times\!10^{-5}$, $2.077\!\times\!10^{-4}$] & 10/10 & 5.76 & 1.88 \\
Pure p=4, d=2 & 256 & $2.494\!\times\!10^{-3}$ & $1.248\!\times\!10^{-2}$ [$1.087\!\times\!10^{-2}$, $1.465\!\times\!10^{-2}$] & 4/10 & 5.01 & 16.12 \\
\bottomrule
\end{tabular}
\end{table}

\FloatBarrier

\subsection{Natural \texorpdfstring{\(V\)}{V}-distance for the \texorpdfstring{\(p\)}{p}-models}

\Cref{fig:baseline-v} repeats the two pure \(p=4\) comparisons with the relative \(V\)-distance defined in
\Cref{sec:numerics}.  Keeping this geometry separate from the \(W^{1,4}\) error is important because the
two quantities need not have the same terminal behavior.  The figure uses the same computed-state and
RFM uncertainty conventions as the Sobolev panels.

\begin{figure}[!htbp]
\centering
\begin{subfigure}[t]{0.49\textwidth}
\includegraphics[width=\linewidth]{C4_v_distance.pdf}
\caption{Pure \(p=4,d=1\).}
\end{subfigure}\hfill
\begin{subfigure}[t]{0.49\textwidth}
\includegraphics[width=\linewidth]{C5_v_distance.pdf}
\caption{Pure \(p=4,d=2\).}
\end{subfigure}
\caption{Relative \(V\)-distance versus effective degrees of freedom.  No theoretical reference line
or common-grid interpolation is drawn.}
\label{fig:baseline-v}
\end{figure}
\FloatBarrier

\subsection{Interpretation limits}

The equal-DOF results compare finite-pool CGA with the specified FEM and RFM implementations.  They do
not establish a runtime ordering, do not compare CGA with every randomized or greedy dictionary method,
and do not turn endpoint ratios into asymptotic complexity estimates. The two-dimensional pure
\(p=4\) RFM summaries use
the four available realizations at the terminal width and are not used as a complete ten-run statistical
comparison. Supplementary energy-gap comparisons are more sensitive to quadrature because they
subtract two separately integrated energies. The supplementary numerical data specify the prescribed and
converged realization counts, interpolation brackets, and computed coefficient counts.
